\section{Charge conventions}\label{app:A.2}

The elementary charge is $e > 0$ with
\begin{equation}
  e^2 = 4\pi\alpha, \qquad \alpha = 1/137.035999,
  \label{eq:A.2-e2}
\end{equation}
the rationalized Heaviside--Lorentz convention of particle physics~\cite{mohr2025codata}. The electron carries charge $-e$ and a nucleus of species $\kappa$ carries $+Z_\kappa e$, and the charge density operator is $e$ times the sum of the nuclear densities weighted by $Z_\kappa$ minus the electron density. Two electrons at distance $r$ have the interaction energy $e^2/(4\pi r) = \alpha/r$, and Gauss's law states that $\nabla\cdot\vect{E}$ equals the charge density, without a factor $4\pi$ (\cref{app:C.9}). The Coulomb interaction in momentum space, with the Fourier convention of \cref{app:A.3}, is the kernel
\begin{equation}
  v_{\mathrm{C}}(k) \equiv \frac{e^2}{k^2} = \frac{4\pi\alpha}{k^2},
  \label{eq:A.2-vC}
\end{equation}
the transform of $e^2/(4\pi r)$ by \cref{eq:C.9-Gq} at $m_V = 0$. In a crystal it is a diagonal matrix in reciprocal-lattice vectors, $(v_{\mathrm{C}})_{\vect{G}\vect{G}'} = v_{\mathrm{C}}(|\vect{q}+\vect{G}|)\,\delta_{\vect{G}\vect{G}'}$ (\cref{sec:4.4}).

\subsection{Heaviside--Lorentz units and the loss function}

Dark-matter couplings are written in field theory, where Heaviside--Lorentz units are standard: the dark-photon coupling to the electron, $y_e = -\zeta e$, is the electromagnetic one scaled by $\zeta$ (\cref{sec:1.4}), and the Yukawa potential of \cref{eq:1.4-yukawa} has the same $1/(4\pi r)$ as Coulomb's law. The dielectric response of a solid is expressed in the same units, so the relation between the structure factor and the loss function carries no factors of $4\pi$. The dynamic structure factor and the loss function are related by
\begin{equation}
  S(q,\omega) = \frac{2q^2}{e^2}\,\frac{L(q,\omega)}{1 - \ee^{-\beta\omega}} = \frac{q^2}{2\pi\alpha}\,\frac{L(q,\omega)}{1 - \ee^{-\beta\omega}} ,
  \label{eq:A.2-SL}
\end{equation}
derived in \cref{app:A.7}~\cites[Eq.~(6)]{knapen2021dark}[Eq.~(S.7)]{hochberg2021determining}. In a dark-photon model the factor $1/e^2$ cancels against $y_e^2 = \zeta^2e^2$, and the rate depends only on $\zeta^2$ and the loss function~\cite[below Eq.~(S.8)]{hochberg2021determining}.

\subsection{Converting from Gaussian and SI units}

Condensed-matter sources usually use Gaussian units, in which the electron charge $e_{\mathrm{G}}$ obeys $e_{\mathrm{G}}^2 = \alpha$ (with $\hbar = c = 1$), the Coulomb energy is $e_{\mathrm{G}}^2/r$ and the Coulomb kernel is $4\pi e_{\mathrm{G}}^2/k^2$. SI sources write $e^2/(4\pi\varepsilon_0^{\mathrm{SI}}r)$, with $\varepsilon_0^{\mathrm{SI}}$ the vacuum permittivity; with $\hbar = c = 1$, $e_{\mathrm{SI}}^2/\varepsilon_0^{\mathrm{SI}} = 4\pi\alpha$. The three charges are therefore related by
\begin{equation}
  e^2 = 4\pi e_{\mathrm{G}}^2 = \frac{e_{\mathrm{SI}}^2}{\varepsilon_0^{\mathrm{SI}}} = 4\pi\alpha .
  \label{eq:A.2-map}
\end{equation}
A formula is converted by replacing every $e^2$ of this book with $4\pi e_{\mathrm{G}}^2$ (to Gaussian) or $e_{\mathrm{SI}}^2/\varepsilon_0^{\mathrm{SI}}$ (to SI), and the reverse for the other direction~\cite[Appendix]{jackson1999classical}. Dimensionless response functions, $\varepsilon$, $\varepsilon^{-1}$ and $L$, have the same value in every system; only the formulas that express them through charges change.

The plasma frequency is
\begin{equation*}
  \omega_p^2 = \frac{e^2\nd_e}{m_e} = \frac{4\pi e_{\mathrm{G}}^2\nd_e}{m_e} = \frac{e_{\mathrm{SI}}^2\nd_e}{\varepsilon_0^{\mathrm{SI}}m_e} ,
\end{equation*}
in the three systems; all three give \SI{16.6}{eV} for the four valence electrons per atom of silicon, $\nd_e = \SI{2.00e23}{cm^{-3}}$ at $a = \SI{5.431}{\angstrom}$ (\cref{sec:4.3}). The dielectric matrix is $\varepsilon_{\vect{G}\vect{G}'} = \delta_{\vect{G}\vect{G}'} - v_{\mathrm{C}}(|\vect{q}+\vect{G}|)\tilde R_{\vect{G}\vect{G}'}$ (\cref{app:A.7}); in the Gaussian sources of the local-field literature~\cite{adler1962quantum,wiser1963dielectric} the same expression appears with $4\pi e_{\mathrm{G}}^2/|\vect{q}+\vect{G}|^2$, and the polarizability there is our $\tilde R$ up to the sign and the factor $1/\Vol$ that each source fixes in its own definition. The optical conductivity enters through $\varepsilon = 1 + \ii\sigma_{\mathrm{cond}}/\omega$, so that $\sigma_1 = \omega\varepsilon_2$; the Gaussian form is $\varepsilon = 1 + 4\pi\ii\sigma_{\mathrm{cond}}/\omega$ and the SI form $\varepsilon = 1 + \ii\sigma_{\mathrm{cond}}/(\varepsilon_0^{\mathrm{SI}}\omega)$, so a measured conductivity $\sigma_1^{\mathrm{SI}}$ in \si{S.m^{-1}} gives $\varepsilon_2 = \sigma_1^{\mathrm{SI}}/(\varepsilon_0^{\mathrm{SI}}\omega)$. The prefactor $8\pi^2\alpha/(\Omega\omega^2)$ of the independent-particle $\varepsilon_2$ in the optical limit (\cref{sec:4.4}) reads $8\pi^2 e_{\mathrm{G}}^2/(\Omega\omega^2)$ in the Gaussian sources, and the two agree under \cref{eq:A.2-map}.

Fields convert as $\vect{E} = \vect{E}_{\mathrm{G}}/\sqrt{4\pi} = (\varepsilon_0^{\mathrm{SI}})^{1/2}\vect{E}_{\mathrm{SI}}$, so that $e\vect{E} = e_{\mathrm{G}}\vect{E}_{\mathrm{G}} = e_{\mathrm{SI}}\vect{E}_{\mathrm{SI}}$ and the force on a charge is the same in every system. Magnetic fields convert in the same way, $\vect{B} = \vect{B}_{\mathrm{G}}/\sqrt{4\pi} = (\varepsilon_0^{\mathrm{SI}})^{1/2}\vect{B}_{\mathrm{SI}}$ with $c = 1$. The Bohr magneton is then $\mu_{\mathrm{B}} = e/(2m_e)$, and the spin coupling of an electron is $g_e\mu_{\mathrm{B}}\,\vect{s}\cdot\vect{B}$, with $\vect{s}$ its spin and $g_e \simeq 2$; the product of $\mu_{\mathrm{B}}$ and the field strength is an energy with the same value in every system, $\SI{5.788e-5}{eV}$ per tesla.

The symbol $\varepsilon_0$ alone is the static dielectric constant of a crystal, lattice included (\cref{sec:5.5}); the vacuum permittivity appears only as $\varepsilon_0^{\mathrm{SI}}$, inside quoted SI formulas.
